% Stand delineation and the HBAD models — formulation and inputs (equations as implemented; code is the reference where docs differ).
% Build: scripts/build.sh   (Tectonic; theme and Fira Sans copied from the PERSEUS deck)
\documentclass[aspectratio=169,11pt]{beamer}
\usepackage{fontspec}
\setsansfont{FiraSans}[Path=fonts/, Extension=.otf, Scale=0.95, UprightFont=*-Regular, BoldFont=*-SemiBold, ItalicFont=*-Italic, BoldItalicFont=*-SemiBold]
\setmonofont{FiraMono}[Extension=.otf, Scale=0.85, UprightFont=*-Regular, BoldFont=*-Bold]
\usepackage{amsmath,amssymb}
\usepackage{unicode-math}
\setmathfont{FiraMath-Regular.otf}
\usepackage{booktabs,tabularx,array}
\usetheme{perseus}
\newcommand{\hl}[1]{{\color{pGreen}\bfseries #1}}
\newcommand{\muted}[1]{{\color{pMuted}#1}}
\newcommand{\src}[1]{\par\vfill{\tiny\color{pMuted}#1\par}}
\setlength{\abovedisplayskip}{3pt}\setlength{\belowdisplayskip}{3pt}
\title{Stand delineation, HBAD \& generative vegetation}
\subtitle{Formulation and inputs — equations as implemented}
\author{BahaaEddin AlAila}
\institute{University of Georgia}
\date{October 2026}
\begin{document}

\begin{frame}
\titlepage
\end{frame}

% ----------------------------------------------------------------------------------------------------------------------------
\begin{frame}{Stand delineation (GA + FL) — inputs and land status}
\footnotesize
\begin{columns}[T]
\begin{column}{0.47\textwidth}
\hl{Inputs} (all on the TreeMap 30\,m EPSG:5070 grid; 1\,px = 0.09\,ha)
\begin{itemize}\setlength\itemsep{1pt}
\item TreeMap 2022 plot ids + attributes (BA, height, QMD, cover, TPA, softwood share)
\item Hansen GFC v1.12: loss year, tree cover 2000, land mask
\item LANDFIRE EVT 2022: excluded types; tree / shrub / herb life form
\item AlphaEarth 2022, 2025: 64-d embeddings, de-quantised $\operatorname{sign}(q)(q/127.5)^2$
\item SOLUS 0–30\,cm: clay, sand, SOC, AWC, pH, restrictive depth
\item Barriers (1-px lines): TIGER roads and water, FL parcels
\end{itemize}
\end{column}
\begin{column}{0.51\textwidth}
\hl{Status classes} (recent = Hansen loss $\ge 2017$; former forest = cover$_{2000}\ge 25\%$)
\begin{itemize}\setlength\itemsep{1pt}
\item \textbf{1 forest}: any TreeMap pixel
\item \textbf{2 harvested}: recent loss on former forest, EVT not excluded (overrides 1)
\item \textbf{3 regenerating}: not TreeMap, not excluded, and older loss, or a disturbed / plantation EVT, or shrub/herb with cover$_{2000}\ge 50\%$
\item \textbf{4 forest unmapped}: EVT tree life form, not TreeMap
\item \textbf{9 converted}: loss on former forest, now an excluded EVT (not forest land)
\item 0 outside the AOI and on Hansen non-land
\end{itemize}
Speckle cleaning: same-class patches $<$ 2\,ha take the class of the nearest larger forest-land patch ($\le$\,15\,px).
\end{column}
\end{columns}
\src{scripts 10\_forest\_land.py, 11\_land\_clean.py, config/pilot.yml}
\end{frame}

% ----------------------------------------------------------------------------------------------------------------------------
\begin{frame}{Stand delineation — the per-pixel feature vector (34 dimensions)}
\footnotesize
Each group is projected and scaled so that its total variance on a fit sample (60 tiles $\times$ 4000 forest-land pixels) equals $w^2$:
\begin{align*}
\mathbf{y}_{\text{AEF}}  &= (\mathbf a_{25}-\boldsymbol\mu)\,V_{16}^{\mathsf{T}}\;1.0\,/\sqrt{\textstyle\sum_{k\le16}\lambda_k} &&\text{16 PCs, AlphaEarth 2025}\\
\mathbf{y}_{\text{old}}  &= (\mathbf a_{22}-\boldsymbol\mu)\,V_{8}^{\mathsf{T}}\;0.4\,/\sqrt{\textstyle\sum_{k\le8}\lambda_k} &&\text{AlphaEarth 2022 on the 2025 basis (8)}\\
y_{\text{chg}} &= \bigl(1-\cos(\mathbf a_{25},\mathbf a_{22})\bigr)\;0.5\,/\operatorname{sd}(1-\cos) &&\text{change since 2022 (1)}\\
\mathbf{y}_{\text{TM}} &= (\mathbf x-\boldsymbol\mu)/\boldsymbol\sigma\;\cdot\;0.6/\sqrt6 &&\text{TreeMap attributes (6)}\\
\mathbf{y}_{\text{soil}} &= \mathbf z\,V_{3}^{\mathsf{T}}\;0.3\,/\sqrt{\textstyle\sum_{k\le3}\lambda_k} &&\text{3 PCs of the z-scored soils}
\end{align*}
\begin{itemize}\setlength\itemsep{1pt}
\item TreeMap: ln BA, height, ln QMD, cover, ln TPA, softwood share; zeroed (= sample mean) off TreeMap forest or where harvested.
\item The weights are per-group standard deviations: AlphaEarth dominates, TreeMap and soils refine.
\end{itemize}
\src{scripts/features.py (Featurizer.features), 20\_fit\_features.py}
\end{frame}

% ----------------------------------------------------------------------------------------------------------------------------
\begin{frame}{Stand delineation — superpixels and the region-adjacency merge}
\footnotesize
\begin{columns}[T]
\begin{column}{0.48\textwidth}
\hl{1. Superpixels} — SLIC per 512-px tile and per status, on forest land minus barrier pixels: $n_{\text{seg}} = \lfloor n_{\text{px}}/11\rfloor$ (1\,ha), connectivity enforced; never cross barriers, statuses or tiles. \hl{Input: raw AlphaEarth 2022 $\oplus$ 2025}, compactness 0.15 (next slide); the merge uses the 34-d stack.

\medskip
\hl{2. Merge} on the region-adjacency graph (edges = shared boundary $L_{ab}$):
\begin{align*}
c_{ab} &= \underbrace{\frac{n_an_b}{n_a+n_b}\,\lVert\boldsymbol\mu_a-\boldsymbol\mu_b\rVert^2}_{\text{Ward}}\times s_{ab}\\
s_{ab} &= 1+0.5\Bigl(1-\min\bigl(1,\tfrac{L_{ab}}{\min(P_a,P_b)}\bigr)\Bigr)
\end{align*}
($s_{ab}$ favours long shared boundaries; $P$ = perimeter)
same status, $n_a+n_b\le$ 200\,ha; each round mutually-best pairs with $c_{ab}<\tau$ merge.
\end{column}
\begin{column}{0.50\textwidth}
\hl{3. Threshold $\tau$} — bisected (14 steps in $[10^{-2},10^{4}]$, 300 sample components) to an area-weighted median stand of 15\,ha.

\medskip
\hl{4. Minimum mapping unit} — then regions $<$ 22\,px (2\,ha) merge into their cheapest neighbour of the same status.

\medskip
\hl{5. Finalise} — barrier and sliver pixels are grown into 4-adjacent stands (same status first), so stands stay connected; fragments $<$ MMU join the same-status neighbour with the longest shared boundary, only toward larger ones (acyclic).

\medskip
\hl{Outputs} — \texttt{stands.tif}, polygons: status, area, loss year, TreeMap means, main plot, soils.
\end{column}
\end{columns}
\src{scripts 21\_superpixels.py, 22\_merge.py, rag\_merge.py, 23\_finalize.py}
\end{frame}

% ----------------------------------------------------------------------------------------------------------------------------
\begin{frame}{Stand delineation — one tunable, verified pipeline and what tuning found}
\scriptsize
\begin{columns}[T]
\begin{column}{0.44\textwidth}
\hl{Pipeline} \texttt{stands\_pipeline.py}: land, clean, fit, (embed), superpixels, merge, finalise, verify, metrics; one config + \texttt{--set} overrides; content-hashed stage cache; 10 invariants checked per run.

\medskip
\hl{Score} on layers the segmentation does not use (canopy 2000, EVC cover, slope, TPI): the share of what \emph{equal-count square blocks} leave unexplained that the stands capture,
\[
\text{share}=\frac{R^2_{\text{stands}}-R^2_{\text{blocks}}}{1-R^2_{\text{blocks}}}
\]
ranked on Baker Co.\ FL, \hl{reported on a second window} (Waycross GA). Oracle (superpixels merged on the layer itself): the stands reach 62\,\% of its gain on canopy 2000.
\end{column}
\begin{column}{0.54\textwidth}
\renewcommand{\arraystretch}{1.05}
\begin{tabular}{@{}lcc@{}}
\toprule
SLIC input, compactness & Baker & Waycr.\\
\midrule
\textbf{AlphaEarth 22+25 (128-d), 0.15} & \textbf{0.197} & \textbf{0.194}\\
IWAE 90\,m latent (12-d), 0.08 & 0.194 & 0.191\\
34-d stack, fixed scaling, 0.15 & 0.182 & 0.178\\
34-d stack, 0.3 (old default) & 0.173 & 0.169\\
\bottomrule
\end{tabular}

\medskip
\begin{itemize}\setlength\itemsep{1pt}
\item merging in the embedding spaces hurt; TreeMap smoothing, fixed scaling: no effect
\item IWAE: conv.\ encoder of 3$\times$3 px, $\mathcal N(0,I)$ prior, heteroscedastic decoder — ties raw AlphaEarth with 10$\times$ fewer channels
\item in progress: finer superpixels, SLIC as the 2nd level, nested stand levels
\end{itemize}
\end{column}
\end{columns}
\src{docs/STANDS\_PIPELINE.md; config/tune\_slic.yml, tune\_embed.yml, tune\_embed2.yml; scripts/stands\_metrics.py, embed\_iwae.py}
\end{frame}

% ----------------------------------------------------------------------------------------------------------------------------
\begin{frame}{HBAD v1 — empirical cohort model (FIA regressions + structural competition)}
\footnotesize
\textbf{State} per cohort: species, age, TPH, QMD $D$ (cm), height $H$; site: $\text{BA}=\tfrac{\pi}{4}(D/100)^2\,\text{TPH}$, $\text{SDI}=\sum\text{TPH}\,(D/25.4)^{1.605}$, $\text{RD}=\text{SDI}/\text{SDI}_{\max}$, BAL = BA of larger cohorts.
\begin{align*}
\Delta D &= G_{\text{mult}}\exp\!\bigl(g_0+g_1\ln\tfrac D{20}+g_2\tfrac D{20}+g_3\ln\tfrac{A}{30}-g_k\tfrac{\text{BAL}}{10}\bigr)\in[0,3]\ \text{cm/yr}\\
H' &= 1.3+(H-1.3)\,\frac{\text{CR}(D')}{\text{CR}(D)},\qquad \text{CR}(D)=h_a\bigl(1-e^{-h_bD}\bigr)^{h_c}\\
h &= M_{\text{mult}}\exp\!\bigl(m_0+m_1\ln\tfrac D{20}+m_2\tfrac D{20}+m_3\text{RD}_c+m_4\tfrac{\text{BAL}}{10}+m_5\text{RD}_c\ln\tfrac D{20}\bigr)\le0.5\\
\text{TPH} &\leftarrow \text{TPH}\,e^{-h},\qquad \text{RD}_c=\min(\text{RD},1.5)
\end{align*}
\vspace{-4pt}
\begin{itemize}\setlength\itemsep{1pt}
\item Reineke self-thinning: while SDI $>$ SDI$_{\max}$, remove trees from the smallest cohorts up; cohorts end at longevity.
\item Establishment: Biomass-Succession block (PROB\_ESTAB; light class from RD; new cohort 750 TPH, 1\,cm).
\item \hl{Fitted}: $G_{\text{mult}}$, $M_{\text{mult}}$ (ecoregion $\times$ species), SDI$_{\max}$ (species). \hl{Fixed}: the regression coefficients (FIA GA+FL).
\item \hl{Inputs}: FIA cohorts (age, AGB $\rightarrow$ TPH, QMD; AGB via $\ln m = a+b\ln(D/2.54)$); no climate, CO$_2$ or soil drivers.
\end{itemize}
\src{landis\_fe docs/hbad\_plugin.md, src/plugins/HBAD/plugin.jl}
\end{frame}

% ----------------------------------------------------------------------------------------------------------------------------
\begin{frame}{HBAD v2 — biomass-free process cohorts with light competition by height}
\footnotesize
\textbf{State} per cohort: $N$ (trees/ha), $D$, $H$, age $A$, vigour quantile $z$. Overtopping by taller cohorts (relative density units):
\[
O_i=\frac{\sum_{H_j>H_i}\text{SDI}_j+\tfrac12\sum_{H_j=H_i}\text{SDI}_j+\tfrac12\text{SDI}_i}{\text{SDI}_{\max}}
\]
\vspace{-6pt}
\begin{align*}
\Delta D &= r_D\,(D/20)^{c_D}\,e^{-D/D_{\max}}\;e^{-g_{OD}O}\,e^{g_{RD}\text{RD}}\,e^{\sigma_v z}\\
\Delta H &= r_H\,(H/15)^{c_H}\,e^{-H/H_{\max}}\;e^{-g_{OH}O+g_{RH}\text{RD}}\,e^{\sigma_v z}\\
h &= f_M\Bigl[e^{m_{bg}+m_D\ln\frac D{20}+m_{D2}\frac D{20}}+m_{\text{sup}}\bigl(1-\min(1,e^{-g_{OD}O}e^{\sigma_vz})\bigr)\bigl(\tfrac D{20}\bigr)^{m_{SD}}\Bigr]+e^{m_A(A/L-1)}\\
N &\leftarrow N\,e^{-h}\\
R &= 185.2\,\exp\bigl(\rho_0+\rho_S\ln(1+\text{BA}_{sp,\,D\ge12.7})+\rho_L\,\text{SDI}/600\bigr)\ \text{trees/ha/yr}
\end{align*}
\vspace{-4pt}
\begin{itemize}\setlength\itemsep{1pt}
\item \hl{Parameters}: $\approx$\,14 per species key ($H_{\max}, r_H, c_H, D_{\max}, r_D, c_D, g_{OD}, g_{RD}, g_{OH}, g_{RH}, m_{bg}, m_D, m_{\text{sup}}$, SDI$_{\max}$), globals, recruitment: $\approx$\,175–210 in all. Recruits enter at 2.54\,cm, 3\,m.
\item \hl{Inputs}: remeasured FIA trees (undisturbed intervals), cohorts by species $\times$ 10-yr birth class $\times$ size class; \hl{fit}: binomial survival + Gaussian $\ln D$, $\ln H$ of survivors.
\end{itemize}
\src{landis\_hbad2 docs/hbad\_v2.md, src/plugins/HBAD2/HBAD2.jl}
\end{frame}

% ----------------------------------------------------------------------------------------------------------------------------
\begin{frame}{HBAD3P — process and trait-bound growth (stem-volume balance)}
\footnotesize
Same cohort state as v2 + crown ratio CR; $O$, RD as in v2. Water: $W$ = growing-season AET/PET (monthly bucket: PRISM precipitation, Thornthwaite PET, plot AWC); $W_n$ = its plot normal.
\begin{align*}
\Delta V &= \varepsilon_k\,(D/20)^{b}\cdot 122.98\;\underbrace{\bigl[f_0+(1-f_0)e^{-aO}\bigr]}_{\text{light}}\;e^{g_R\text{RD}}\;\underbrace{W^{\omega}}_{\text{water}}\;e^{\sigma_v z}\;\Bigl(\frac{\text{CR}}{0.4}\Bigr)^{\varphi}\\
&\qquad\times\underbrace{\max\bigl(0,1-(H/1.1H_{\max})^{\theta}\bigr)}_{\text{hydraulic}}\;-\;\underbrace{\rho\,D^2H}_{\text{maintenance}}\\
\Delta D &= \frac{\Delta V}{2DH}\\
H^*(D) &= 1.37+(H_{\max}-1.37)\bigl(1-e^{-sD/(H_{\max}-1.37)}\bigr)^{c}\\
\Delta H &= \bigl[H^*(D+\Delta D)-H^*(D)\bigr]\bigl(1+\eta(1-L_h)\bigr)+\lambda\bigl(H^*(D)(1+\mu(1-L_h))-H\bigr),\quad L_h=e^{-k_hO}\\
\text{CR}^* &= \text{cr}_{lo}+(\text{cr}_{hi}-\text{cr}_{lo})\,e^{-k_cO},\qquad \text{CR}\rightarrow\text{CR}^*\ \text{at rates }\tau_r\ (\text{recession}),\ \tau_e\ (\text{expansion}),\ \text{CR}_0=\text{FIA CR}
\end{align*}
\vspace{-4pt}
\begin{itemize}\setlength\itemsep{1pt}
\item Per species \hl{traits fixed} ($H_{\max}$, shade tolerance ST, longevity $L$); per key only the growth efficiency $\varepsilon_k$ and fecundity $F_k$ are fitted. 122.98 is a fixed unit constant of $\varepsilon$.
\item Densities per ha of the \emph{forested} subplot; one pipeline \texttt{pipeline3p.sh <REGION>} with 37 checks.
\end{itemize}
\src{landis\_hbad3 docs/hbad3\_design.md \S4, src/plugins/HBAD3/HBAD3P.jl (step\_site!)}
\end{frame}

% ----------------------------------------------------------------------------------------------------------------------------
\begin{frame}{HBAD3P — mortality, thinning, recruitment, climate}
\footnotesize
\begin{align*}
g_m &= e^{-a_mO}\,e^{g_R\text{RD}}\,W^{\omega}\,e^{\sigma_vz}\,(\text{CR}/0.4)^{\varphi}\qquad\text{(survival carbon balance)}\\
h &= f_M\Bigl[\frac{m_0}{L}+m_s\,e^{\sigma_{ST}(\text{ST}-3)}\,e^{-\beta g_m}\Bigl(\frac{D}{20}\Bigr)^{\zeta}\Bigr]+e^{m_A(A/L-1)}\\
f_M &= e^{p_M[\text{planted}]}\;e^{\,m_W(W_n-W)/0.1}\qquad\text{(drought anomaly)}\\
N &\leftarrow N\,e^{-h} \qquad\text{and, while RD}>1\text{: remove SDI}\propto\text{SDI}_i(1-v_i)\\
R_k &= F_k\;\exp\bigl(-\lambda_0e^{\lambda_1(\text{ST}-3)}\,\text{RD}\bigr)\times(\text{seed source}\,?\,1:\varepsilon_0)\\
&\qquad\text{fitted on sapling ingrowth}\times\text{probe-recruit survival to the next visit}
\end{align*}
\begin{columns}[T]
\begin{column}{0.48\textwidth}
\hl{Climate inputs}
\begin{itemize}\setlength\itemsep{1pt}
\item growth: $W^{\omega}$ (water limitation)
\item mortality: drought anomaly $m_W(W_n-W)$
\item CO$_2$, temperature: scenario assumptions from the literature (not identifiable from FIA)
\end{itemize}
\end{column}
\begin{column}{0.50\textwidth}
\hl{Parameters}: $\approx$\,30–33 per region (vs $\approx$\,175–210 in v2):
\begin{itemize}\setlength\itemsep{0pt}
\item per key: $\varepsilon_k$, $F_k$
\item growth: $b$, $a$, $g_R$, $\theta$, $\omega$, $a_m$, $\rho$, $f_0$, $\varphi$
\item height: $s$, $c$, $\eta$, $\lambda$, $\mu$, $k_h$
\item mortality: $m_0$, $m_s$, $\sigma_{ST}$, $\beta$, $\zeta$, $m_W$
\item planted $p_D$, $p_H$, $p_M$; vigour $\sigma_v$; crown; recruitment
\end{itemize}
\end{column}
\end{columns}
\src{src/plugins/HBAD3/HBAD3P.jl; docs/hbad3\_design.md \S4, \S7}
\end{frame}

% ----------------------------------------------------------------------------------------------------------------------------
\begin{frame}{The three HBAD generations side by side}
\scriptsize
\renewcommand{\arraystretch}{1.1}
\begin{tabularx}{\textwidth}{>{\bfseries}l X X X}
\toprule
 & HBAD v1 & HBAD v2 & HBAD3P \\
\midrule
growth & FIA $\Delta D$ regression $\times G_{\text{mult}}$, BAL & power law in $D$, $H$ $\times$ light ($O$), crowding (RD), vigour & stem-volume balance: light, crowding, water $W^{\omega}$, crown, hydraulic limit \\
height & Chapman–Richards ratio & own growth equation & trait $H$–$D$ curve + light rule \\
mortality / thinning & regression hazard $\times M_{\text{mult}}$ (RD, BAL); Reineke cap & background + suppression + senescence; SDI cap & carbon-balance hazard + drought anomaly + senescence; RD $\le1$ \\
recruitment & Biomass-Succession block & stand-table regression & fecundity $\times$ light window $\times$ seed source \\
climate & none & none & $W$ (PRISM + AWC), drought anomaly \\
parameters & $G,M$ per eco$\times$spp, SDI$_{\max}$ & $\approx$\,175–210 & $\approx$\,30–33 per region \\
fit (GA+FL val NLL) & — & $-272{,}655$ & $-273{,}469$ \\
\bottomrule
\end{tabularx}
\src{fit: Sim A likelihood (binomial survival + Gaussian ln D, ln H), CMA-ES. HBAD3P free runs (val, pipeline 2026-10-22): BA +4.5 / +0.5\,\%, QMD +5.4 / +8.3\,\%, Lorey H +3.2 / +2.0\,\% (GA+FL / APP); stems $\ge$12.7\,cm $-13$ / $-19$\,\% (in-growth too low)}
\end{frame}

% ----------------------------------------------------------------------------------------------------------------------------
\begin{frame}{Generative vegetation — FIA plots without their true locations}
\footnotesize
\begin{columns}[T]
\begin{column}{0.48\textwidth}
\hl{Goal}: stand-level vegetation (cohorts: species, age, $N$, QMD, $H$, CR) mapped from rasters, \emph{independent of TreeMap}, learned from FIA.

\medskip
\hl{Training table} (GA+FL, 8{,}712 plots, 34{,}300 visits since 2000): cohort tokens per visit (species $\times$ age class: $\ln N$, $\ln$QMD, $H$, CR) per forested area; consecutive visits = pairs.

\medskip
\hl{The location is latent}: public coordinates are fuzzed within $\approx$1.6\,km. Candidates: 90\,m footprints of forest land in the disk, $\le$256 per plot (stratified, inclusion-weighted), with AlphaEarth 2022 / 2025, disturbance history 2000–2025 (LANDFIRE class, severity; Hansen loss), soils.
\end{column}
\begin{column}{0.50\textwidth}
\hl{Location prior from FIA's own records}: TRT / DSTRB codes matched \emph{by type} to the footprint's evidence in the event year:
\begin{itemize}\setlength\itemsep{0pt}
\item cutting / site prep.\ / regeneration $\leftrightarrow$ harvest, thinning, chemical, fire, \emph{unknown}, Hansen loss: LR 4
\item crown fire: 4; ground fire, insects, weather, human: 2
\item records of one event (within 2\,yr) count once
\item no code over a remeasurement $\Rightarrow$ stand-replacing evidence in that period $\times0.25$
\end{itemize}
\hl{History consistency}: footprints compatible with \emph{all} detectable records, in time order (unsatisfiable ones skipped).

\medskip
\muted{Location information per plot is tiny ($\approx$0.003 nats): the records anchor the \emph{dynamics}, not the place.}
\end{column}
\end{columns}
\src{g02b\_build\_gafl.py (EVK), g02e\_fia\_records.py, g02g\_history\_consistency.py; location tests g02f, g08}
\end{frame}

% ----------------------------------------------------------------------------------------------------------------------------
\begin{frame}{Generative vegetation — temporal JEPA with disturbance jumps}
\footnotesize
\begin{align*}
z_v &= E(\text{tokens}_v)\in\mathbb R^{16} &&\text{masked set encoder, SIGReg toward }\mathcal N(0,I)\\
\tfrac{dz}{dt} &= -\bigl(\tfrac1\tau+g\bigr)\odot z+g\odot A,\quad g=\sigma(W[z,s]+b) &&\text{liquid time-constant flow (yearly), site }s\\
z &\leftarrow z+J\bigl([z,\ \text{class},\ \text{severity},\ \text{Hansen}]\bigr) &&\text{jump at a candidate's disturbance years}\\
\mathcal L_{\text{MIL}} &= -\log\tfrac1K\textstyle\sum_{c=1}^{K}e^{-d_c/T},\quad c\sim\text{prior (iid)} &&d_c=\lVert h_{22}(z_{2022})-a^c_{22}\rVert^2+\lVert h_{25}(z_{2025})-a^c_{25}\rVert^2,\ K{=}48
\end{align*}
\vspace{-6pt}
\scriptsize
\begin{columns}[T]
\begin{column}{0.5\textwidth}
\begin{itemize}\setlength\itemsep{1pt}
\item \hl{pairs}: $z(v_1)$ flowed along FIA-consistent candidate histories to $v_2$ $\approx z(v_2)$ and $v_2$'s cohorts: the records anchor the jumps
\item \hl{decoder}: cohort slots (presence + $\ln N$, $\ln$QMD, $H$, CR); summary loss TPH / BA / QMD + Bray–Curtis
\end{itemize}
\end{column}
\begin{column}{0.48\textwidth}
\begin{itemize}\setlength\itemsep{1pt}
\item \hl{A}: landscape prior $p(z\mid X_L,\text{year})$, mixture-density net (8 Gaussians) on the whole 1.6\,km disk — robust to the fuzzing
\item \hl{B}: likelihood $p(\text{AEF}\mid z,\text{history})$, fine-tuned on well-localised plots
\end{itemize}
\end{column}
\end{columns}
\src{scripts g03b\_train\_gafl.py (E, flow, jumps, heads, MIL, pairs, decoder), g04\_bayes\_downscale.py (A, B, calibration)}
\end{frame}

% ----------------------------------------------------------------------------------------------------------------------------
\begin{frame}{Generative vegetation — mapping stands, verification, results}
\footnotesize
\begin{columns}[T]
\begin{column}{0.46\textwidth}
\hl{Mapping} (per stand of the delineation): $X_L$ of the nearest 1\,km landscape node (same descriptor as the plots' disks) $\rightarrow$ $z\sim$ A at 2021 $\rightarrow$ optionally weighted by B on the stand's footprints $\rightarrow$ flow to 2025 through the footprints' disturbance histories $\rightarrow$ decode; calibration factors from held-out calibration plots (15\,\% of training blocks).

\medskip
\hl{Pipeline} \texttt{genveg\_pipeline.py}: table $\rightarrow$ records $\rightarrow$ history $\rightarrow$ stand footprints $\rightarrow$ train (2 seeds) $\rightarrow$ dynamics test $\rightarrow$ downscale $\rightarrow$ map $\rightarrow$ Forest Explorer $\rightarrow$ verify: \hl{45 independent checks} (SQL / direct raster reads / re-implemented specification) pass.
\end{column}
\begin{column}{0.52\textwidth}
\hl{Held-out visits 2022–2025, mapping path} (median $|\log\text{error}|$; spatial blocks, 3.2\,km buffer)

\smallskip
\renewcommand{\arraystretch}{1.05}
\begin{tabular}{@{}lccc@{}}
\toprule
 & TPH & BA & QMD\\
\midrule
generative (calibrated) & 0.611 & \textbf{0.412} & \textbf{0.308}\\
kNN on landscape features & 0.664 & 0.414 & 0.359\\
null (median) & \textbf{0.609} & 0.453 & 0.327\\
\bottomrule
\end{tabular}

\medskip
\begin{itemize}\setlength\itemsep{1pt}
\item modest: beats kNN on QMD, null on BA; no better than null on TPH
\item dominant-species accuracy 0.42 vs kNN 0.36 (majority 0.22)
\item open: the direct predictor's BA level varies with the training seed (log bias $-0.27\ldots-0.62$; calibration absorbs it)
\end{itemize}
\end{column}
\end{columns}
\src{scripts g05\_map\_stands.py, gv\_verify.py; docs/GENVEG\_PIPELINE.md, GENERATIVE\_VEG\_PLAN.md}
\end{frame}

\end{document}
